% ============================================================
% Dodatki do _formatting.cls dla estetycznych tabel wynikowych
% Wklej po:
% \RequirePackage{multirow}
% \RequirePackage{tabularx,array}
% ============================================================

\RequirePackage{makecell}  % łamanie nagłówków tabel: \thead{...\\...}
\RequirePackage{ragged2e}  % estetyczne wyrównanie tekstu w kolumnach X
\RequirePackage{siunitx}   % wyrównywanie liczb w kolumnach S

\sisetup{
  detect-all,
  output-decimal-marker = {,},
  group-separator = {\,},
  group-minimum-digits = 4,
  table-number-alignment = center
}

% Kolumny pomocnicze do tabularx.
% Y: tekst zawijany do szerokości tabeli, wyrównany do lewej.
% L/C: kolumny o stałej szerokości, jeśli będą potrzebne w innych rozdziałach.
\newcolumntype{Y}{>{\RaggedRight\arraybackslash}X}
\newcolumntype{L}[1]{>{\RaggedRight\arraybackslash}p{#1}}
\newcolumntype{C}[1]{>{\Centering\arraybackslash}p{#1}}

% Spójny wygląd nagłówków i odstępów w tabelach.
\renewcommand\theadfont{\bfseries\small}
\renewcommand\theadalign{cc}
\setlength{\tabcolsep}{3.5pt}
\setlength{\arrayrulewidth}{0.4pt}
\renewcommand{\arraystretch}{1.2}
